\documentclass[11pt]{article}

\usepackage[utf8]{inputenc}
\usepackage[T1]{fontenc}
\usepackage{lmodern}
\usepackage{geometry}
\usepackage{amsmath,amssymb,amsfonts,amsthm,mathtools}
\usepackage{bm}
\usepackage{enumitem}
\usepackage{microtype}
\usepackage{hyperref}
\usepackage{titlesec}
\usepackage{booktabs}
\usepackage{array}
\usepackage{setspace}
\usepackage{mathrsfs}
\usepackage{physics}

\geometry{margin=1in}
\hypersetup{colorlinks=true, linkcolor=black, urlcolor=blue, citecolor=black}
\setlist[enumerate]{topsep=0.4em,itemsep=0.35em}
\setlist[itemize]{topsep=0.3em,itemsep=0.25em}
\titleformat{\section}{\large\bfseries}{\thesection.}{0.5em}{}
\titleformat{\subsection}{\normalsize\bfseries}{\thesubsection.}{0.5em}{}
\setstretch{1.06}

\newcommand{\Aether}{\AE{}ther}
\newcommand{\Aflow}{\AE{}ther-flow}
\newcommand{\TheoryName}{\Aflow{} Interpretation of Relativity}
\newcommand{\TheoryProgram}{\Aether{} / \Aflow{} framework}

\newcommand{\CompletedMetric}{g^{\sharp}}
\newcommand{\MatterSpace}{\Psi}

\newcommand{\Car}{\mathrm{car}}
\newcommand{\Cur}{\mathrm{cur}}
\newcommand{\Hom}{\mathrm{hom}}

\newcommand{\ForSpace}[1]{\mathcal I_{#1}^{\mathrm{for}}}
\newcommand{\PrimShell}{\mathcal S_{\mathrm{prim}}}
\newcommand{\MetricOut}{\mathscr G_{\mathrm{out}}}
\newcommand{\MatterOut}{\mathscr T_{\mathrm{out}}}

\newcommand{\SubSectionPackage}{\mathfrak V^{\mathrm{subsec}}}
\newcommand{\MicroSectionPackage}{\mathfrak W^{\mathrm{microsec}}}
\newcommand{\NanoSectionPackage}{\mathfrak X^{\mathrm{nanosec}}}
\newcommand{\PicoSectionPackage}{\mathfrak Y^{\mathrm{picosec}}}
\newcommand{\PicoMinSectionCore}{\mathfrak Y^{\mathrm{pic,sec}}}
\newcommand{\FemtoSectionPackage}{\mathfrak Z^{\mathrm{femtosec}}}
\newcommand{\AttoSectionPackage}{\mathfrak A^{\mathrm{attosec}}}
\newcommand{\ZeptoImageCore}{\mathfrak B^{\mathrm{zep,img}}}
\newcommand{\ZeptoMinSectionCore}{\mathfrak B^{\mathrm{zep,sec}}}
\newcommand{\YoctoImageCore}{\mathfrak C^{\mathrm{yoc,img}}}
\newcommand{\YoctoMinSectionCore}{\mathfrak C^{\mathrm{yoc,sec}}}
\newcommand{\YoctoSelectorPackage}{\mathfrak C^{\mathrm{yoc,sel}}}

\newcommand{\YoctoSelectorSpace}[1]{\mathcal D_{#1}^{\mathrm{yoc,sel}}}
\newcommand{\YoctoSelectorBody}[1]{\mathcal B_{#1}^{\mathrm{yoc,sel}}}
\newcommand{\YoctoSelectorFunctional}[1]{\mathscr L_{#1}^{\mathrm{yoc,sel}}}
\newcommand{\YoctoSelectorShell}[1]{\mathcal Z_{#1}^{\mathrm{yoc,sel}}}

\newcommand{\PreSelectorPackage}{\mathfrak D^{\mathrm{presel}}}
\newcommand{\PreSelectorMinSectionCore}{\mathfrak D^{\mathrm{presel,sec}}}

\newcommand{\PreSectionSelectorPackage}{\mathfrak E^{\mathrm{presel,sel}}}
\newcommand{\PreSectionSelectorSpace}[1]{\mathcal F_{#1}^{\mathrm{presel,sel}}}
\newcommand{\PreSectionSelectorBody}[1]{\mathcal B_{#1}^{\mathrm{presel,sel}}}
\newcommand{\PreSectionSelectorProj}[1]{\Pi_{#1}^{\mathrm{presel,sel}}}
\newcommand{\PreSectionSelectorFiber}[2]{\mathcal G_{#1}^{\mathrm{presel,sel}}(#2)}

\newcommand{\KernelPackage}{\mathfrak F^{\mathrm{presel,ker}}}
\newcommand{\KernelSpace}[1]{\mathcal H_{#1}^{\mathrm{presel,ker}}}
\newcommand{\KernelBody}[1]{\mathcal K_{#1}^{\mathrm{presel,ker}}}
\newcommand{\KernelProj}[1]{\Lambda_{#1}^{\mathrm{presel,ker}}}
\newcommand{\KernelFiber}[2]{\mathcal J_{#1}^{\mathrm{presel,ker}}(#2)}
\newcommand{\KernelComposedFiber}[2]{\mathcal M_{#1}^{\mathrm{presel,ker}}(#2)}

\newcommand{\PreKernelPackage}{\mathfrak G^{\mathrm{presel,preker}}}
\newcommand{\PreKernelSpace}[1]{\mathcal N_{#1}^{\mathrm{presel,preker}}}
\newcommand{\PreKernelBody}[1]{\mathcal P_{#1}^{\mathrm{presel,preker}}}
\newcommand{\PreKernelProj}[1]{\Upsilon_{#1}^{\mathrm{presel,preker}}}
\newcommand{\PreKernelFiber}[2]{\mathcal U_{#1}^{\mathrm{presel,preker}}(#2)}
\newcommand{\PreKernelFullFiber}[2]{\mathcal W_{#1}^{\mathrm{presel,preker}}(#2)}

\newcommand{\ProtoKernelPackage}{\mathfrak H^{\mathrm{presel,proker}}}
\newcommand{\ProtoKernelSpace}[1]{\mathcal R_{#1}^{\mathrm{presel,proker}}}
\newcommand{\ProtoKernelBody}[1]{\mathcal T_{#1}^{\mathrm{presel,proker}}}
\newcommand{\ProtoKernelProj}[1]{\Omega_{#1}^{\mathrm{presel,proker}}}
\newcommand{\ProtoKernelFiber}[2]{\mathcal X_{#1}^{\mathrm{presel,proker}}(#2)}
\newcommand{\ProtoKernelKernelFiber}[2]{\mathcal L_{#1}^{\mathrm{presel,proker}}(#2)}
\newcommand{\ProtoKernelSelectorFiber}[2]{\mathcal A_{#1}^{\mathrm{presel,proker}}(#2)}
\newcommand{\ProtoKernelFullFiber}[2]{\mathcal C_{#1}^{\mathrm{presel,proker}}(#2)}

\newcommand{\PreProtoKernelPackage}{\mathfrak I^{\mathrm{presel,preproker}}}
\newcommand{\PreProtoKernelSpace}[1]{\mathcal S_{#1}^{\mathrm{presel,preproker}}}
\newcommand{\PreProtoKernelBody}[1]{\mathcal D_{#1}^{\mathrm{presel,preproker}}}
\newcommand{\PreProtoKernelProj}[1]{\Xi_{#1}^{\mathrm{presel,preproker}}}
\newcommand{\PreProtoKernelFiber}[2]{\mathcal B_{#1}^{\mathrm{presel,preproker}}(#2)}
\newcommand{\PreProtoKernelFunctional}[1]{\mathscr V_{#1}^{\mathrm{presel,preproker}}}
\newcommand{\PreProtoKernelMinLift}[1]{\mathfrak z_{#1}^{\mathrm{presel,preproker}}}
\newcommand{\PreProtoKernelPreKernelFiber}[2]{\mathcal M_{#1}^{\mathrm{presel,preproker,preker}}(#2)}
\newcommand{\PreProtoKernelPreKernelMin}[1]{\mathfrak a_{#1}^{\mathrm{presel,preproker}}}
\newcommand{\PreProtoKernelKernelFiber}[2]{\mathcal N_{#1}^{\mathrm{presel,preproker,ker}}(#2)}
\newcommand{\PreProtoKernelKernelMin}[1]{\mathfrak b_{#1}^{\mathrm{presel,preproker}}}
\newcommand{\PreProtoKernelSelectorFiber}[2]{\mathcal P_{#1}^{\mathrm{presel,preproker,sel}}(#2)}
\newcommand{\PreProtoKernelSelectorMin}[1]{\mathfrak c_{#1}^{\mathrm{presel,preproker}}}
\newcommand{\PreProtoKernelFullFiber}[2]{\mathcal Q_{#1}^{\mathrm{presel,preproker,full}}(#2)}
\newcommand{\PreProtoKernelFullMin}[1]{\mathfrak d_{#1}^{\mathrm{presel,preproker}}}
\newcommand{\PreProtoKernelShell}[1]{\mathcal O_{#1}^{\mathrm{presel,preproker}}}

\newcommand{\PrePreProtoKernelPackage}{\mathfrak J^{\mathrm{presel,prepreproker}}}
\newcommand{\PrePreProtoKernelSpace}[1]{\mathcal I_{#1}^{\mathrm{presel,prepreproker}}}
\newcommand{\PrePreProtoKernelBody}[1]{\mathcal L_{#1}^{\mathrm{presel,prepreproker}}}
\newcommand{\PrePreProtoKernelProj}[1]{\Theta_{#1}^{\mathrm{presel,prepreproker}}}
\newcommand{\PrePreProtoKernelFiber}[2]{\mathcal R_{#1}^{\mathrm{presel,prepreproker}}(#2)}
\newcommand{\PrePreProtoKernelFunctional}[1]{\mathscr W_{#1}^{\mathrm{presel,prepreproker}}}
\newcommand{\PrePreProtoKernelMinLift}[1]{\mathfrak e_{#1}^{\mathrm{presel,prepreproker}}}
\newcommand{\PrePreProtoKernelProtoFiber}[2]{\mathcal T_{#1}^{\mathrm{presel,prepreproker,proto}}(#2)}
\newcommand{\PrePreProtoKernelProtoMin}[1]{\mathfrak f_{#1}^{\mathrm{presel,prepreproker}}}
\newcommand{\PrePreProtoKernelPreKernelFiber}[2]{\mathcal V_{#1}^{\mathrm{presel,prepreproker,preker}}(#2)}
\newcommand{\PrePreProtoKernelPreKernelMin}[1]{\mathfrak g_{#1}^{\mathrm{presel,prepreproker}}}
\newcommand{\PrePreProtoKernelKernelFiber}[2]{\mathcal W_{#1}^{\mathrm{presel,prepreproker,ker}}(#2)}
\newcommand{\PrePreProtoKernelKernelMin}[1]{\mathfrak h_{#1}^{\mathrm{presel,prepreproker}}}
\newcommand{\PrePreProtoKernelSelectorFiber}[2]{\mathcal Y_{#1}^{\mathrm{presel,prepreproker,sel}}(#2)}
\newcommand{\PrePreProtoKernelSelectorMin}[1]{\mathfrak i_{#1}^{\mathrm{presel,prepreproker}}}
\newcommand{\PrePreProtoKernelFullFiber}[2]{\mathcal Z_{#1}^{\mathrm{presel,prepreproker,full}}(#2)}
\newcommand{\PrePreProtoKernelFullMin}[1]{\mathfrak j_{#1}^{\mathrm{presel,prepreproker}}}
\newcommand{\PrePreProtoKernelShell}[1]{\mathcal K_{#1}^{\mathrm{presel,prepreproker}}}

\newcommand{\InducedSelectorFunctional}[1]{\widetilde{\mathscr Q}_{#1}^{\mathrm{presel,sel}}}
\newcommand{\InducedSelectorShell}[1]{\widetilde{\mathcal Z}_{#1}^{\mathrm{presel,sel}}}
\newcommand{\InducedKernelFunctional}[1]{\widetilde{\mathscr R}_{#1}^{\mathrm{presel,ker}}}
\newcommand{\InducedKernelMinLift}[1]{\widetilde{\mathfrak r}_{#1}^{\mathrm{presel,ker}}}
\newcommand{\InducedKernelComposedMin}[1]{\widetilde{\mathfrak m}_{#1}^{\mathrm{presel,ker}}}
\newcommand{\InducedKernelShell}[1]{\widetilde{\mathcal Y}_{#1}^{\mathrm{presel,ker}}}
\newcommand{\InducedPreKernelFunctional}[1]{\widetilde{\mathscr S}_{#1}^{\mathrm{presel,preker}}}
\newcommand{\InducedPreKernelMinLift}[1]{\widetilde{\mathfrak s}_{#1}^{\mathrm{presel,preker}}}
\newcommand{\InducedPreKernelShell}[1]{\widetilde{\mathcal O}_{#1}^{\mathrm{presel,preker}}}
\newcommand{\InducedProtoKernelFunctional}[1]{\widetilde{\mathscr U}_{#1}^{\mathrm{presel,proker}}}
\newcommand{\InducedProtoKernelMinLift}[1]{\widetilde{\mathfrak v}_{#1}^{\mathrm{presel,proker}}}
\newcommand{\InducedProtoKernelKernelMin}[1]{\widetilde{\mathfrak w}_{#1}^{\mathrm{presel,proker}}}
\newcommand{\InducedProtoKernelSelectorMin}[1]{\widetilde{\mathfrak x}_{#1}^{\mathrm{presel,proker}}}
\newcommand{\InducedProtoKernelFullMin}[1]{\widetilde{\mathfrak y}_{#1}^{\mathrm{presel,proker}}}
\newcommand{\InducedProtoKernelShell}[1]{\widetilde{\mathcal E}_{#1}^{\mathrm{presel,proker}}}
\newcommand{\InducedPreProtoKernelFunctional}[1]{\widetilde{\mathscr V}_{#1}^{\mathrm{presel,preproker}}}
\newcommand{\InducedPreProtoKernelMinLift}[1]{\widetilde{\mathfrak z}_{#1}^{\mathrm{presel,preproker}}}
\newcommand{\InducedPreProtoKernelPreKernelMin}[1]{\widetilde{\mathfrak a}_{#1}^{\mathrm{presel,preproker}}}
\newcommand{\InducedPreProtoKernelKernelMin}[1]{\widetilde{\mathfrak b}_{#1}^{\mathrm{presel,preproker}}}
\newcommand{\InducedPreProtoKernelSelectorMin}[1]{\widetilde{\mathfrak c}_{#1}^{\mathrm{presel,preproker}}}
\newcommand{\InducedPreProtoKernelFullMin}[1]{\widetilde{\mathfrak d}_{#1}^{\mathrm{presel,preproker}}}
\newcommand{\InducedPreProtoKernelShell}[1]{\widetilde{\mathcal O}_{#1}^{\mathrm{presel,preproker}}}

\newtheorem{definition}{Definition}
\newtheorem{proposition}{Proposition}
\newtheorem{remark}{Remark}
\newtheorem{corollary}{Corollary}
\newtheorem{theorem}{Theorem}

\title{\textbf{The \TheoryName{}}\\[0.4em]
\large Primitive-Reservoir Observer-Reduction\\
\large Proto-Kernel-Dynamics-Generating Substrate\\
\large Pre-Selector Pre-Proto-Kernel Dynamics\\
\large from Pre-Proto-Kernel-Dynamics-Generating\\
\large Substrate Pre-Selector Pre-Pre-Proto-Kernel Dynamics}
\author{}
\date{}

\begin{document}

\maketitle

\begin{abstract}
The current primitive-reservoir observer-reduction frontier already sharpened the live burden below the current pre-selector proto-kernel line: the current proto-kernel-dynamics-generating substrate pre-selector pre-proto-kernel dynamics package \(\PreProtoKernelPackage\) is now the live unexplained datum, while the current proto-kernel package \(\ProtoKernelPackage\), the current pre-kernel package \(\PreKernelPackage\), the current kernel package \(\KernelPackage\), and all downstream observer-relevant outputs have already been spent. The active routing note therefore demands the next honest move below that frontier: either derive or justify \(\PreProtoKernelPackage\) from still more primitive substrate structure in a way that genuinely changes benchmark-facing status, or prove a sharper theorem-level collapse strictly inside it.

The present note takes the default first route. For each sector it introduces pre-proto-kernel-dynamics-generating substrate pre-selector pre-pre-proto-kernel dynamics \(\PrePreProtoKernelPackage\): a compact convex pre-pre-proto-kernel body over the already fixed current pre-proto-kernel-body geometry, an affine pre-pre-proto-kernel projection to that geometry, a nonnegative smooth pre-pre-proto-kernel functional with unique inner fiber minimizers, a proto-kernel-scale minimizing law forcing the current pre-proto-kernel minimizer lift, a pre-kernel-scale minimizing law forcing the current pre-proto-kernel pre-kernel minimizer, a current-kernel-scale minimizing law forcing the current pre-proto-kernel current-kernel minimizer, a selector-scale minimizing law forcing the current pre-proto-kernel selector minimizer, and a benchmark-shell-scale minimizing shell whose projection fixes the current exact pre-proto-kernel shell. The current pre-proto-kernel functional is then no longer an independent primitive stopping point. It is the effective minimum value of the deeper pre-pre-proto-kernel law on the inner pre-pre-proto-kernel fibers, the current exact pre-proto-kernel shell is the projected exact pre-pre-proto-kernel shell over the already fixed selector benchmark shell, and the current pre-proto-kernel proto-scale, pre-kernel-scale, current-kernel-scale, selector-scale, and benchmark-shell-scale minimizers are all outer projections of unique deeper pre-pre-proto-kernel minimizers.

The benchmark boundary remains fixed throughout. The one operative metric is still exactly \(\CompletedMetric\), the ordinary matter route is unchanged, there is no second operative metric, the low-energy matter law is not enlarged, and no stronger promotion language is licensed. The positive result is therefore a genuine benchmark-facing gain below the current pre-selector pre-proto-kernel frontier: the remaining honest burden is no longer the current pre-selector pre-proto-kernel dynamics as such, but the deeper pre-selector pre-pre-proto-kernel dynamics that canonically generate them, or a still smaller collapse strictly inside that deeper package.
\end{abstract}

\section{Introduction}

The active derivational program of \TheoryName{} again reaches the point at which depth alone is not enough. The current active-primary theorem already showed that the current pre-selector pre-proto-kernel package is not the primitive stopping point at present scope. Once the proto-kernel-dynamics-generating substrate pre-selector pre-proto-kernel dynamics package \(\PreProtoKernelPackage\) is fixed, the current pre-proto-kernel functional, the current pre-proto-kernel minimizer lift, the current pre-proto-kernel pre-kernel minimizer, the current pre-proto-kernel current-kernel minimizer, the current pre-proto-kernel selector minimizer, the current exact pre-proto-kernel shell, and the current benchmark-shell-scale full pre-proto-kernel minimizer are forced canonically from it. The routing consequence is immediate. If the project goal is to move the same benchmark one-metric package toward a disciplined first-principles recovery, the next honest benchmark-facing manuscript must either derive or justify \(\PreProtoKernelPackage\) itself from still more primitive substrate structure, or prove a genuinely smaller theorem-level collapse strictly inside it.

The present note takes the direct derivation route. It does not claim that no sharper internal collapse will ever be available. It records only that no smaller observer-relevant datum has yet been fixed on the active line strictly inside the current pre-proto-kernel package without already presupposing that package's law. The disciplined next move is therefore to spend \(\PreProtoKernelPackage\) itself. Concretely, the note introduces a deeper pre-pre-proto-kernel layer below the current pre-proto-kernel package and proves that the current pre-proto-kernel functional, the current pre-proto-kernel minimizer hierarchy, and the current exact pre-proto-kernel shell are all canonical outputs of that deeper pre-pre-proto-kernel law.

\paragraph{Framework and claim boundary.}
\TheoryProgram{} denotes the broader \Aether{} / \Aflow{} framework built on the claims that the \Aether{} is the underlying four-dimensional substrate of reality, the \Aflow{} is its intrinsic ordered motion, observed three-dimensional space is the local experiential slice of that deeper substrate, S-time is the experienced order of change arising from matter, light, and the \Aflow{}, and observed expansion is the three-dimensional appearance of deeper four-dimensional ordered motion. Within that framework, \TheoryName{} denotes the exact-closure theory in which the effective gravitational dynamics are adopted to be exactly Einsteinian with universal matter coupling. In the active sequence, \emph{adoption} means use of that established relativistic dynamics without claiming substrate derivation, while \emph{derivation} is reserved for a first-principles recovery from explicit substrate variables.

The present note stays strictly inside that boundary. It does not alter the benchmark exact-closure package. It does not introduce a second operative metric or a new low-energy matter law. It only strengthens the deeper origin-side explanation of why the current pre-selector pre-proto-kernel package used by the immediately preceding theorem should itself hold.

\section{Fixed Inputs}

\begin{definition}[Recorded primitive continuation data]
\label{def:fixed}
Fix the following continuation data from the active primitive-reservoir observer-reduction line.
\begin{enumerate}
    \item For each sector label \(\sigma\in\{\Car,\Cur,\Hom\}\), a primitive forcing shell
    \[
    \ForSpace{\sigma}.
    \]
    \item The product shell
    \[
    \PrimShell
    \equiv
    \ForSpace{\Car}\times \ForSpace{\Cur}\times \ForSpace{\Hom}.
    \]
    \item The recorded metric and ordinary-matter outputs
    \[
    \MetricOut:\PrimShell\to\{\text{observer metrics}\},
    \qquad
    \MatterOut:\PrimShell\times\MatterSpace\to\{\text{observer stress tensors}\},
    \]
    satisfying
    \begin{equation}
    \MetricOut(\mathbf w)=\CompletedMetric,
    \qquad
    \MatterOut(\mathbf w,\psi)
    =
    T_{\mu\nu}^{\mathrm{matter}}[\CompletedMetric,\psi]
    \label{eq:fixedoutputs}
    \end{equation}
    for every \(\mathbf w\in\PrimShell\) and every matter configuration \(\psi\in\MatterSpace\).
\end{enumerate}
\end{definition}

\begin{definition}[Recorded selector benchmark base]
\label{def:selectorbase}
Fix \(\sigma\in\{\Car,\Cur,\Hom\}\). The active line already fixes the following benchmark-facing selector data:
\begin{enumerate}
    \item a finite-dimensional Euclidean selector state space
    \[
    \YoctoSelectorSpace{\sigma};
    \]
    \item a compact convex selector body
    \[
    \YoctoSelectorBody{\sigma}\subset\YoctoSelectorSpace{\sigma}
    \]
    with nonempty interior;
    \item a nonnegative smooth selector functional
    \[
    \YoctoSelectorFunctional{\sigma}:\YoctoSelectorSpace{\sigma}\to[0,\infty);
    \]
    \item a closed embedded exact selector shell
    \[
    \YoctoSelectorShell{\sigma}\subset\YoctoSelectorBody{\sigma};
    \]
    \item benchmark faithfulness, meaning preservation of the benchmark outputs \eqref{eq:fixedoutputs}, no second operative metric, no enlarged low-energy matter law, and use of this selector benchmark base as already fixed upstream data on the active line.
\end{enumerate}
\end{definition}

\begin{definition}[Recorded current selector-body geometry]
\label{def:selectorgeometry}
Fix \(\sigma\in\{\Car,\Cur,\Hom\}\). The active line already fixes the following current selector-body geometry:
\begin{enumerate}
    \item a finite-dimensional Euclidean current selector state space
    \[
    \PreSectionSelectorSpace{\sigma};
    \]
    \item a compact convex current selector body
    \[
    \PreSectionSelectorBody{\sigma}\subset\PreSectionSelectorSpace{\sigma}
    \]
    with nonempty interior;
    \item an affine surjection
    \[
    \PreSectionSelectorProj{\sigma}:\PreSectionSelectorSpace{\sigma}\to\YoctoSelectorBody{\sigma}
    \]
    satisfying
    \[
    \PreSectionSelectorProj{\sigma}\bigl(\PreSectionSelectorBody{\sigma}\bigr)=\YoctoSelectorBody{\sigma};
    \]
    \item for each \(y\in\YoctoSelectorBody{\sigma}\), the current selector fiber
    \[
    \PreSectionSelectorFiber{\sigma}{y}
    \equiv
    \bigl(\PreSectionSelectorProj{\sigma}\bigr)^{-1}(y)\cap\PreSectionSelectorBody{\sigma},
    \]
    which is required to be nonempty, compact, and convex.
\end{enumerate}
\end{definition}

\begin{definition}[Recorded current kernel-body geometry]
\label{def:kernelgeometry}
Fix \(\sigma\in\{\Car,\Cur,\Hom\}\). The active line already fixes the following current kernel-body geometry:
\begin{enumerate}
    \item a finite-dimensional Euclidean current kernel state space
    \[
    \KernelSpace{\sigma};
    \]
    \item a compact convex current kernel body
    \[
    \KernelBody{\sigma}\subset\KernelSpace{\sigma}
    \]
    with nonempty interior;
    \item an affine surjection
    \[
    \KernelProj{\sigma}:\KernelSpace{\sigma}\to\PreSectionSelectorSpace{\sigma}
    \]
    satisfying
    \[
    \KernelProj{\sigma}\bigl(\KernelBody{\sigma}\bigr)=\PreSectionSelectorBody{\sigma};
    \]
    \item for each \(u\in\PreSectionSelectorBody{\sigma}\), the current kernel fiber
    \[
    \KernelFiber{\sigma}{u}
    \equiv
    \bigl(\KernelProj{\sigma}\bigr)^{-1}(u)\cap\KernelBody{\sigma},
    \]
    which is required to be nonempty, compact, and convex.
\end{enumerate}
\end{definition}

\begin{definition}[Recorded current pre-kernel-body geometry]
\label{def:prekernelgeometry}
Fix \(\sigma\in\{\Car,\Cur,\Hom\}\). The active line already fixes the following current pre-kernel-body geometry:
\begin{enumerate}
    \item a finite-dimensional Euclidean current pre-kernel state space
    \[
    \PreKernelSpace{\sigma};
    \]
    \item a compact convex current pre-kernel body
    \[
    \PreKernelBody{\sigma}\subset\PreKernelSpace{\sigma}
    \]
    with nonempty interior;
    \item an affine surjection
    \[
    \PreKernelProj{\sigma}:\PreKernelSpace{\sigma}\to\KernelSpace{\sigma}
    \]
    satisfying
    \[
    \PreKernelProj{\sigma}\bigl(\PreKernelBody{\sigma}\bigr)=\KernelBody{\sigma};
    \]
    \item for each \(w\in\KernelBody{\sigma}\), the current pre-kernel fiber
    \[
    \PreKernelFiber{\sigma}{w}
    \equiv
    \bigl(\PreKernelProj{\sigma}\bigr)^{-1}(w)\cap\PreKernelBody{\sigma},
    \]
    which is required to be nonempty, compact, and convex.
\end{enumerate}
\end{definition}

\begin{definition}[Recorded current proto-kernel-body geometry]
\label{def:protokernelgeometry}
Fix \(\sigma\in\{\Car,\Cur,\Hom\}\). The active line already fixes the following current proto-kernel-body geometry:
\begin{enumerate}
    \item a finite-dimensional Euclidean current proto-kernel state space
    \[
    \ProtoKernelSpace{\sigma};
    \]
    \item a compact convex current proto-kernel body
    \[
    \ProtoKernelBody{\sigma}\subset\ProtoKernelSpace{\sigma}
    \]
    with nonempty interior;
    \item an affine surjection
    \[
    \ProtoKernelProj{\sigma}:\ProtoKernelSpace{\sigma}\to\PreKernelSpace{\sigma}
    \]
    satisfying
    \[
    \ProtoKernelProj{\sigma}\bigl(\ProtoKernelBody{\sigma}\bigr)=\PreKernelBody{\sigma};
    \]
    \item for each \(q\in\PreKernelBody{\sigma}\), the current proto-kernel fiber
    \[
    \ProtoKernelFiber{\sigma}{q}
    \equiv
    \bigl(\ProtoKernelProj{\sigma}\bigr)^{-1}(q)\cap\ProtoKernelBody{\sigma},
    \]
    which is required to be nonempty, compact, and convex.
\end{enumerate}
\end{definition}

\begin{definition}[Recorded current pre-proto-kernel-body geometry]
\label{def:preprotogeometry}
Fix \(\sigma\in\{\Car,\Cur,\Hom\}\). The active line already fixes the following current pre-proto-kernel-body geometry below the current proto-kernel-body geometry:
\begin{enumerate}
    \item a finite-dimensional Euclidean current pre-proto-kernel state space
    \[
    \PreProtoKernelSpace{\sigma};
    \]
    \item a compact convex current pre-proto-kernel body
    \[
    \PreProtoKernelBody{\sigma}\subset\PreProtoKernelSpace{\sigma}
    \]
    with nonempty interior;
    \item an affine surjection
    \[
    \PreProtoKernelProj{\sigma}:\PreProtoKernelSpace{\sigma}\to\ProtoKernelSpace{\sigma}
    \]
    satisfying
    \[
    \PreProtoKernelProj{\sigma}\bigl(\PreProtoKernelBody{\sigma}\bigr)=\ProtoKernelBody{\sigma};
    \]
    \item for each \(z\in\ProtoKernelBody{\sigma}\), the current pre-proto-kernel fiber
    \[
    \PreProtoKernelFiber{\sigma}{z}
    \equiv
    \bigl(\PreProtoKernelProj{\sigma}\bigr)^{-1}(z)\cap\PreProtoKernelBody{\sigma},
    \]
    which is required to be nonempty, compact, and convex.
\end{enumerate}
\end{definition}

\begin{definition}[Current pre-selector pre-proto-kernel-dynamics target]
\label{def:preprototarget}
Fix \(\sigma\in\{\Car,\Cur,\Hom\}\) and the recorded current selector-body, kernel-body, pre-kernel-body, proto-kernel-body, and pre-proto-kernel-body geometries of Definitions \ref{def:selectorgeometry}--\ref{def:preprotogeometry}. A \emph{benchmark-faithful proto-kernel-dynamics-generating substrate pre-selector pre-proto-kernel dynamics package} \(\PreProtoKernelPackage\) on that geometry consists of the following data.
\begin{enumerate}
    \item A nonnegative smooth pre-proto-kernel functional
    \[
    \PreProtoKernelFunctional{\sigma}:\PreProtoKernelSpace{\sigma}\to[0,\infty)
    \]
    such that, for every \(z\in\ProtoKernelBody{\sigma}\), its restriction to \(\PreProtoKernelFiber{\sigma}{z}\) is strictly convex and attains a unique minimizer depending smoothly on \(z\); write
    \[
    \PreProtoKernelMinLift{\sigma}:\ProtoKernelBody{\sigma}\to\PreProtoKernelBody{\sigma}.
    \]
    The induced current proto-kernel functional
    \[
    \InducedProtoKernelFunctional{\sigma}:\ProtoKernelSpace{\sigma}\to[0,\infty)
    \]
    is required to be smooth and to satisfy
    \[
    \InducedProtoKernelFunctional{\sigma}(z)
    \equiv
    \PreProtoKernelFunctional{\sigma}\bigl(\PreProtoKernelMinLift{\sigma}(z)\bigr)
    =
    \min_{r\in\PreProtoKernelFiber{\sigma}{z}}
    \PreProtoKernelFunctional{\sigma}(r)
    \qquad
    \text{for every }
    z\in\ProtoKernelBody{\sigma}.
    \]
    Moreover, for every \(q\in\PreKernelBody{\sigma}\), the restriction of \(\InducedProtoKernelFunctional{\sigma}\) to \(\ProtoKernelFiber{\sigma}{q}\) is strictly convex and attains a unique minimizer depending smoothly on \(q\); write
    \[
    \InducedProtoKernelMinLift{\sigma}:\PreKernelBody{\sigma}\to\ProtoKernelBody{\sigma}.
    \]
    \item For each \(q\in\PreKernelBody{\sigma}\), define the current pre-kernel-scale composed pre-proto-kernel fiber
    \[
    \PreProtoKernelPreKernelFiber{\sigma}{q}
    \equiv
    \left\{
    r\in\PreProtoKernelBody{\sigma}:
    \ProtoKernelProj{\sigma}\bigl(\PreProtoKernelProj{\sigma}(r)\bigr)=q
    \right\}.
    \]
    The restriction of \(\PreProtoKernelFunctional{\sigma}\) to \(\PreProtoKernelPreKernelFiber{\sigma}{q}\) is required to attain a unique minimizer depending smoothly on \(q\); write
    \[
    \PreProtoKernelPreKernelMin{\sigma}:\PreKernelBody{\sigma}\to\PreProtoKernelBody{\sigma}.
    \]
    Define the induced current pre-kernel functional
    \[
    \InducedPreKernelFunctional{\sigma}:\PreKernelSpace{\sigma}\to[0,\infty)
    \]
    by
    \[
    \InducedPreKernelFunctional{\sigma}(q)
    \equiv
    \InducedProtoKernelFunctional{\sigma}\bigl(\InducedProtoKernelMinLift{\sigma}(q)\bigr)
    =
    \min_{z\in\ProtoKernelFiber{\sigma}{q}}
    \InducedProtoKernelFunctional{\sigma}(z)
    \qquad
    \text{for every }
    q\in\PreKernelBody{\sigma}.
    \]
    This function is required to be smooth, and for every \(w\in\KernelBody{\sigma}\), its restriction to \(\PreKernelFiber{\sigma}{w}\) is strictly convex and attains a unique minimizer depending smoothly on \(w\); write
    \[
    \InducedPreKernelMinLift{\sigma}:\KernelBody{\sigma}\to\PreKernelBody{\sigma}.
    \]
    \item For each \(w\in\KernelBody{\sigma}\), define the current-kernel-scale composed pre-proto-kernel fiber
    \[
    \PreProtoKernelKernelFiber{\sigma}{w}
    \equiv
    \left\{
    r\in\PreProtoKernelBody{\sigma}:
    \PreKernelProj{\sigma}\bigl(\ProtoKernelProj{\sigma}(\PreProtoKernelProj{\sigma}(r))\bigr)=w
    \right\},
    \]
    and require the restriction of \(\PreProtoKernelFunctional{\sigma}\) to that fiber to attain a unique minimizer depending smoothly on \(w\); write
    \[
    \PreProtoKernelKernelMin{\sigma}:\KernelBody{\sigma}\to\PreProtoKernelBody{\sigma}.
    \]
    Define the induced current kernel functional
    \[
    \InducedKernelFunctional{\sigma}:\KernelSpace{\sigma}\to[0,\infty)
    \]
    by
    \[
    \InducedKernelFunctional{\sigma}(w)
    \equiv
    \InducedPreKernelFunctional{\sigma}\bigl(\InducedPreKernelMinLift{\sigma}(w)\bigr)
    =
    \min_{q\in\PreKernelFiber{\sigma}{w}}
    \InducedPreKernelFunctional{\sigma}(q)
    \qquad
    \text{for every }
    w\in\KernelBody{\sigma}.
    \]
    This function is required to be smooth, and for every \(u\in\PreSectionSelectorBody{\sigma}\), its restriction to \(\KernelFiber{\sigma}{u}\) is strictly convex and attains a unique minimizer depending smoothly on \(u\); write
    \[
    \InducedKernelMinLift{\sigma}:\PreSectionSelectorBody{\sigma}\to\KernelBody{\sigma}.
    \]
    \item For each \(u\in\PreSectionSelectorBody{\sigma}\), define the selector-scale composed pre-proto-kernel fiber
    \[
    \PreProtoKernelSelectorFiber{\sigma}{u}
    \equiv
    \left\{
    r\in\PreProtoKernelBody{\sigma}:
    \KernelProj{\sigma}\bigl(\PreKernelProj{\sigma}(\ProtoKernelProj{\sigma}(\PreProtoKernelProj{\sigma}(r)))\bigr)=u
    \right\},
    \]
    and require the restriction of \(\PreProtoKernelFunctional{\sigma}\) to that fiber to attain a unique minimizer depending smoothly on \(u\); write
    \[
    \PreProtoKernelSelectorMin{\sigma}:\PreSectionSelectorBody{\sigma}\to\PreProtoKernelBody{\sigma}.
    \]
    Define the induced outer selector functional
    \[
    \InducedSelectorFunctional{\sigma}:\PreSectionSelectorSpace{\sigma}\to[0,\infty)
    \]
    by
    \[
    \InducedSelectorFunctional{\sigma}(u)
    \equiv
    \InducedKernelFunctional{\sigma}\bigl(\InducedKernelMinLift{\sigma}(u)\bigr)
    =
    \min_{w\in\KernelFiber{\sigma}{u}}
    \InducedKernelFunctional{\sigma}(w)
    \qquad
    \text{for every }
    u\in\PreSectionSelectorBody{\sigma}.
    \]
    This function is required to be smooth, and for every \(y\in\YoctoSelectorBody{\sigma}\), its restriction to \(\PreSectionSelectorFiber{\sigma}{y}\) is required to be strictly convex and to attain a unique minimizer depending smoothly on \(y\).
    \item For each \(y\in\YoctoSelectorBody{\sigma}\), the restriction of \(\InducedKernelFunctional{\sigma}\) to
    \[
    \KernelComposedFiber{\sigma}{y}
    \equiv
    \left\{
    w\in\KernelBody{\sigma}:
    \PreSectionSelectorProj{\sigma}\bigl(\KernelProj{\sigma}(w)\bigr)=y
    \right\}
    \]
    is required to attain a unique minimizer depending smoothly on \(y\); write
    \[
    \InducedKernelComposedMin{\sigma}:\YoctoSelectorBody{\sigma}\to\KernelBody{\sigma}.
    \]
    Define the induced exact selector shell
    \[
    \InducedSelectorShell{\sigma}
    \equiv
    \left\{
    u\in\PreSectionSelectorBody{\sigma}:
    \PreSectionSelectorProj{\sigma}(u)\in\YoctoSelectorShell{\sigma},
    \;
    \InducedSelectorFunctional{\sigma}(u)
    =
    \min_{v\in\PreSectionSelectorFiber{\sigma}{\PreSectionSelectorProj{\sigma}(u)}}
    \InducedSelectorFunctional{\sigma}(v)
    \right\},
    \]
    the induced exact kernel shell
    \[
    \InducedKernelShell{\sigma}
    \equiv
    \left\{
    w\in\KernelBody{\sigma}:
    \PreSectionSelectorProj{\sigma}\bigl(\KernelProj{\sigma}(w)\bigr)\in\YoctoSelectorShell{\sigma},
    \;
    \InducedKernelFunctional{\sigma}(w)
    =
    \min_{q\in\KernelComposedFiber{\sigma}{\PreSectionSelectorProj{\sigma}(\KernelProj{\sigma}(w))}}
    \InducedKernelFunctional{\sigma}(q)
    \right\},
    \]
    and the induced exact pre-kernel shell
    \[
    \InducedPreKernelShell{\sigma}
    \equiv
    \left\{
    q\in\PreKernelBody{\sigma}:
    \PreSectionSelectorProj{\sigma}\bigl(\KernelProj{\sigma}(\PreKernelProj{\sigma}(q))\bigr)\in\YoctoSelectorShell{\sigma},
    \;
    \InducedPreKernelFunctional{\sigma}(q)
    =
    \min_{p\in\PreKernelFullFiber{\sigma}{\PreSectionSelectorProj{\sigma}(\KernelProj{\sigma}(\PreKernelProj{\sigma}(q)))}} \InducedPreKernelFunctional{\sigma}(p)
    \right\}.
    \]
    The shells \(\InducedKernelShell{\sigma}\) and \(\InducedPreKernelShell{\sigma}\) are required to be closed embedded submanifolds, and
    \[
    \KernelProj{\sigma}\big|_{\InducedKernelShell{\sigma}}
    :
    \InducedKernelShell{\sigma}\xrightarrow{\cong}\InducedSelectorShell{\sigma},
    \qquad
    \PreKernelProj{\sigma}\big|_{\InducedPreKernelShell{\sigma}}
    :
    \InducedPreKernelShell{\sigma}\xrightarrow{\cong}\InducedKernelShell{\sigma}
    \]
    are required to be diffeomorphisms.
    \item For each \(y\in\YoctoSelectorBody{\sigma}\), require the restriction of \(\InducedProtoKernelFunctional{\sigma}\) to
    \[
    \ProtoKernelFullFiber{\sigma}{y}
    \equiv
    \left\{
    z\in\ProtoKernelBody{\sigma}:
    \PreSectionSelectorProj{\sigma}\bigl(\KernelProj{\sigma}(\PreKernelProj{\sigma}(\ProtoKernelProj{\sigma}(z)))\bigr)=y
    \right\}
    \]
    to attain a unique minimizer depending smoothly on \(y\); write
    \[
    \InducedProtoKernelFullMin{\sigma}:\YoctoSelectorBody{\sigma}\to\ProtoKernelBody{\sigma}.
    \]
    Define the induced exact proto-kernel shell
    \[
    \InducedProtoKernelShell{\sigma}
    \equiv
    \left\{
    z\in\ProtoKernelBody{\sigma}:
    \PreSectionSelectorProj{\sigma}\bigl(\KernelProj{\sigma}(\PreKernelProj{\sigma}(\ProtoKernelProj{\sigma}(z)))\bigr)\in\YoctoSelectorShell{\sigma},
    \;
    \InducedProtoKernelFunctional{\sigma}(z)
    =
    \min_{p\in\ProtoKernelFullFiber{\sigma}{\PreSectionSelectorProj{\sigma}(\KernelProj{\sigma}(\PreKernelProj{\sigma}(\ProtoKernelProj{\sigma}(z))))}}
    \InducedProtoKernelFunctional{\sigma}(p)
    \right\}.
    \]
    The shell \(\InducedProtoKernelShell{\sigma}\) is required to be a closed embedded submanifold, and
    \[
    \ProtoKernelProj{\sigma}\big|_{\InducedProtoKernelShell{\sigma}}
    :
    \InducedProtoKernelShell{\sigma}\xrightarrow{\cong}\InducedPreKernelShell{\sigma}
    \]
    is required to be a diffeomorphism.
    \item For each \(y\in\YoctoSelectorBody{\sigma}\), define the full pre-proto-kernel fiber
    \[
    \PreProtoKernelFullFiber{\sigma}{y}
    \equiv
    \left\{
    r\in\PreProtoKernelBody{\sigma}:
    \PreSectionSelectorProj{\sigma}\bigl(\KernelProj{\sigma}(\PreKernelProj{\sigma}(\ProtoKernelProj{\sigma}(\PreProtoKernelProj{\sigma}(r))))\bigr)=y
    \right\}.
    \]
    The restriction of \(\PreProtoKernelFunctional{\sigma}\) to \(\PreProtoKernelFullFiber{\sigma}{y}\) is required to attain a unique minimizer depending smoothly on \(y\); write
    \[
    \PreProtoKernelFullMin{\sigma}:\YoctoSelectorBody{\sigma}\to\PreProtoKernelBody{\sigma}.
    \]
    The exact pre-proto-kernel shell
    \[
    \PreProtoKernelShell{\sigma}
    \equiv
    \left\{
    r\in\PreProtoKernelBody{\sigma}:
    \PreSectionSelectorProj{\sigma}\bigl(\KernelProj{\sigma}(\PreKernelProj{\sigma}(\ProtoKernelProj{\sigma}(\PreProtoKernelProj{\sigma}(r))))\bigr)\in\YoctoSelectorShell{\sigma},
    \;
    \PreProtoKernelFunctional{\sigma}(r)
    =
    \min_{s\in\PreProtoKernelFullFiber{\sigma}{\PreSectionSelectorProj{\sigma}(\KernelProj{\sigma}(\PreKernelProj{\sigma}(\ProtoKernelProj{\sigma}(\PreProtoKernelProj{\sigma}(r)))))}}
    \PreProtoKernelFunctional{\sigma}(s)
    \right\}
    \]
    is required to be a closed embedded submanifold, and
    \[
    \PreProtoKernelProj{\sigma}\big|_{\PreProtoKernelShell{\sigma}}
    :
    \PreProtoKernelShell{\sigma}\xrightarrow{\cong}\InducedProtoKernelShell{\sigma}
    \]
    is required to be a diffeomorphism.
    \item Benchmark faithfulness, meaning preservation of the benchmark outputs \eqref{eq:fixedoutputs}, no second operative metric, no enlarged low-energy matter law, and presentation as the current substrate-side pre-proto-kernel-dynamics package rather than as a replacement benchmark theory.
\end{enumerate}
\end{definition}

\begin{definition}[Recorded downstream chain below the current pre-proto-kernel-dynamics frontier]
\label{def:downstream}
Fix \(\sigma\in\{\Car,\Cur,\Hom\}\). The active line already fixes the following downstream data and consequences:
\begin{enumerate}
    \item the current proto-kernel package \(\ProtoKernelPackage\), the current pre-kernel package \(\PreKernelPackage\), the current kernel package \(\KernelPackage\), the current pre-selector-minimizer-section-generating substrate selector dynamics package \(\PreSectionSelectorPackage\), the current observer-relevant pre-selector minimizer-section core \(\PreSelectorMinSectionCore\), the current selector package \(\YoctoSelectorPackage\), the current observer-relevant yocto-section minimizer-section core \(\YoctoMinSectionCore\), the current observer-relevant yocto-section minimizer-image core \(\YoctoImageCore\), the current observer-relevant zepto-section minimizer-section core \(\ZeptoMinSectionCore\), the current observer-relevant zepto-section minimizer-image core \(\ZeptoImageCore\), the current atto-section package \(\AttoSectionPackage\), the current femto-section package \(\FemtoSectionPackage\), the current pico-section package \(\PicoSectionPackage\), the current observer-relevant pico-section minimizer-section core \(\PicoMinSectionCore\), the current nano-section package \(\NanoSectionPackage\), the current micro-section package \(\MicroSectionPackage\), the current sub-section package \(\SubSectionPackage\), and the previously recorded shell-transport consequences used on the active line;
    \item the previously recorded fact that, once a benchmark-faithful current pre-proto-kernel package in the sense of Definition \ref{def:preprototarget} is available on the recorded geometry, the current proto-kernel package is recovered, and all of the downstream data listed in item (1) follow unchanged.
\end{enumerate}
\end{definition}

\begin{remark}
Definitions \ref{def:fixed}--\ref{def:downstream} fix the current frontier package and its benchmark boundary. The one operative metric, the ordinary matter route, the selector benchmark base, the current selector-body geometry, the current kernel-body geometry, the current pre-kernel-body geometry, the current proto-kernel-body geometry, the current pre-proto-kernel-body geometry, and the already recorded downstream consequences are not reopened here. The present note asks only why the current pre-proto-kernel package itself should arise from still more primitive substrate structure.
\end{remark}

\section{Pre-Proto-Kernel-Dynamics-Generating Substrate Pre-Selector Pre-Pre-Proto-Kernel Dynamics}

\begin{definition}[Pre-proto-kernel-dynamics-generating substrate pre-selector pre-pre-proto-kernel dynamics]
\label{def:prepreproto}
Fix \(\sigma\in\{\Car,\Cur,\Hom\}\) and the recorded geometries of Definitions \ref{def:selectorgeometry}--\ref{def:preprotogeometry}. A \emph{pre-proto-kernel-dynamics-generating substrate pre-selector pre-pre-proto-kernel dynamics package} \(\PrePreProtoKernelPackage\) for sector \(\sigma\) consists of the following data.
\begin{enumerate}
    \item A finite-dimensional Euclidean pre-pre-proto-kernel state space \(\PrePreProtoKernelSpace{\sigma}\), a compact convex pre-pre-proto-kernel body
    \[
    \PrePreProtoKernelBody{\sigma}\subset\PrePreProtoKernelSpace{\sigma}
    \]
    with nonempty interior, and an affine surjection
    \[
    \PrePreProtoKernelProj{\sigma}:\PrePreProtoKernelSpace{\sigma}\to\PreProtoKernelSpace{\sigma}
    \]
    satisfying
    \[
    \PrePreProtoKernelProj{\sigma}\bigl(\PrePreProtoKernelBody{\sigma}\bigr)=\PreProtoKernelBody{\sigma}.
    \]
    For each \(r\in\PreProtoKernelBody{\sigma}\), write
    \[
    \PrePreProtoKernelFiber{\sigma}{r}
    \equiv
    \bigl(\PrePreProtoKernelProj{\sigma}\bigr)^{-1}(r)\cap\PrePreProtoKernelBody{\sigma}.
    \]
    Each such fiber is required to be nonempty, compact, and convex.
    \item A nonnegative smooth pre-pre-proto-kernel functional
    \[
    \PrePreProtoKernelFunctional{\sigma}:\PrePreProtoKernelSpace{\sigma}\to[0,\infty)
    \]
    such that, for every \(r\in\PreProtoKernelBody{\sigma}\), its restriction to \(\PrePreProtoKernelFiber{\sigma}{r}\) is strictly convex and attains a unique minimizer depending smoothly on \(r\); write
    \[
    \PrePreProtoKernelMinLift{\sigma}:\PreProtoKernelBody{\sigma}\to\PrePreProtoKernelBody{\sigma}.
    \]
    The induced current pre-proto-kernel functional
    \[
    \InducedPreProtoKernelFunctional{\sigma}:\PreProtoKernelSpace{\sigma}\to[0,\infty)
    \]
    is required to be smooth and to satisfy
    \[
    \InducedPreProtoKernelFunctional{\sigma}(r)
    \equiv
    \PrePreProtoKernelFunctional{\sigma}\bigl(\PrePreProtoKernelMinLift{\sigma}(r)\bigr)
    =
    \min_{s\in\PrePreProtoKernelFiber{\sigma}{r}}
    \PrePreProtoKernelFunctional{\sigma}(s)
    \qquad
    \text{for every }
    r\in\PreProtoKernelBody{\sigma}.
    \]
    Moreover, for every \(z\in\ProtoKernelBody{\sigma}\), the restriction of \(\InducedPreProtoKernelFunctional{\sigma}\) to \(\PreProtoKernelFiber{\sigma}{z}\) is strictly convex and attains a unique minimizer depending smoothly on \(z\); write
    \[
    \InducedPreProtoKernelMinLift{\sigma}:\ProtoKernelBody{\sigma}\to\PreProtoKernelBody{\sigma}.
    \]
    \item For each \(z\in\ProtoKernelBody{\sigma}\), define the current proto-kernel-scale composed pre-pre-proto-kernel fiber
    \[
    \PrePreProtoKernelProtoFiber{\sigma}{z}
    \equiv
    \left\{
    s\in\PrePreProtoKernelBody{\sigma}:
    \PreProtoKernelProj{\sigma}\bigl(\PrePreProtoKernelProj{\sigma}(s)\bigr)=z
    \right\}.
    \]
    The restriction of \(\PrePreProtoKernelFunctional{\sigma}\) to \(\PrePreProtoKernelProtoFiber{\sigma}{z}\) is required to attain a unique minimizer depending smoothly on \(z\); write
    \[
    \PrePreProtoKernelProtoMin{\sigma}:\ProtoKernelBody{\sigma}\to\PrePreProtoKernelBody{\sigma}.
    \]
    Define the induced current proto-kernel functional
    \[
    \InducedProtoKernelFunctional{\sigma}:\ProtoKernelSpace{\sigma}\to[0,\infty)
    \]
    by
    \[
    \InducedProtoKernelFunctional{\sigma}(z)
    \equiv
    \InducedPreProtoKernelFunctional{\sigma}\bigl(\InducedPreProtoKernelMinLift{\sigma}(z)\bigr)
    =
    \min_{r\in\PreProtoKernelFiber{\sigma}{z}}
    \InducedPreProtoKernelFunctional{\sigma}(r)
    \qquad
    \text{for every }
    z\in\ProtoKernelBody{\sigma}.
    \]
    This function is required to be smooth, and for every \(q\in\PreKernelBody{\sigma}\), its restriction to \(\ProtoKernelFiber{\sigma}{q}\) is strictly convex and attains a unique minimizer depending smoothly on \(q\); write
    \[
    \InducedProtoKernelMinLift{\sigma}:\PreKernelBody{\sigma}\to\ProtoKernelBody{\sigma}.
    \]
    \item For each \(q\in\PreKernelBody{\sigma}\), define the current pre-kernel-scale composed pre-pre-proto-kernel fiber
    \[
    \PrePreProtoKernelPreKernelFiber{\sigma}{q}
    \equiv
    \left\{
    s\in\PrePreProtoKernelBody{\sigma}:
    \ProtoKernelProj{\sigma}\bigl(\PreProtoKernelProj{\sigma}(\PrePreProtoKernelProj{\sigma}(s))\bigr)=q
    \right\}.
    \]
    The restriction of \(\PrePreProtoKernelFunctional{\sigma}\) to \(\PrePreProtoKernelPreKernelFiber{\sigma}{q}\) is required to attain a unique minimizer depending smoothly on \(q\); write
    \[
    \PrePreProtoKernelPreKernelMin{\sigma}:\PreKernelBody{\sigma}\to\PrePreProtoKernelBody{\sigma}.
    \]
    The restriction of \(\InducedPreProtoKernelFunctional{\sigma}\) to \(\PreProtoKernelPreKernelFiber{\sigma}{q}\) is also required to attain a unique minimizer depending smoothly on \(q\); write
    \[
    \InducedPreProtoKernelPreKernelMin{\sigma}:\PreKernelBody{\sigma}\to\PreProtoKernelBody{\sigma}.
    \]
    Define the induced current pre-kernel functional
    \[
    \InducedPreKernelFunctional{\sigma}:\PreKernelSpace{\sigma}\to[0,\infty)
    \]
    by
    \[
    \InducedPreKernelFunctional{\sigma}(q)
    \equiv
    \InducedProtoKernelFunctional{\sigma}\bigl(\InducedProtoKernelMinLift{\sigma}(q)\bigr)
    =
    \min_{z\in\ProtoKernelFiber{\sigma}{q}}
    \InducedProtoKernelFunctional{\sigma}(z)
    \qquad
    \text{for every }
    q\in\PreKernelBody{\sigma}.
    \]
    This function is required to be smooth, and for every \(w\in\KernelBody{\sigma}\), its restriction to \(\PreKernelFiber{\sigma}{w}\) is strictly convex and attains a unique minimizer depending smoothly on \(w\); write
    \[
    \InducedPreKernelMinLift{\sigma}:\KernelBody{\sigma}\to\PreKernelBody{\sigma}.
    \]
    \item For each \(w\in\KernelBody{\sigma}\), define the current-kernel-scale composed pre-pre-proto-kernel fiber
    \[
    \PrePreProtoKernelKernelFiber{\sigma}{w}
    \equiv
    \left\{
    s\in\PrePreProtoKernelBody{\sigma}:
    \PreKernelProj{\sigma}\bigl(\ProtoKernelProj{\sigma}(\PreProtoKernelProj{\sigma}(\PrePreProtoKernelProj{\sigma}(s)))\bigr)=w
    \right\}.
    \]
    The restriction of \(\PrePreProtoKernelFunctional{\sigma}\) to that fiber is required to attain a unique minimizer depending smoothly on \(w\); write
    \[
    \PrePreProtoKernelKernelMin{\sigma}:\KernelBody{\sigma}\to\PrePreProtoKernelBody{\sigma}.
    \]
    The restrictions of \(\InducedPreProtoKernelFunctional{\sigma}\) and \(\InducedProtoKernelFunctional{\sigma}\) to \(\PreProtoKernelKernelFiber{\sigma}{w}\) and \(\ProtoKernelKernelFiber{\sigma}{w}\), respectively, are also required to attain unique minimizers depending smoothly on \(w\); write
    \[
    \InducedPreProtoKernelKernelMin{\sigma}:\KernelBody{\sigma}\to\PreProtoKernelBody{\sigma},
    \qquad
    \InducedProtoKernelKernelMin{\sigma}:\KernelBody{\sigma}\to\ProtoKernelBody{\sigma}.
    \]
    Define the induced current kernel functional
    \[
    \InducedKernelFunctional{\sigma}:\KernelSpace{\sigma}\to[0,\infty)
    \]
    by
    \[
    \InducedKernelFunctional{\sigma}(w)
    \equiv
    \InducedPreKernelFunctional{\sigma}\bigl(\InducedPreKernelMinLift{\sigma}(w)\bigr)
    =
    \min_{q\in\PreKernelFiber{\sigma}{w}}
    \InducedPreKernelFunctional{\sigma}(q)
    \qquad
    \text{for every }
    w\in\KernelBody{\sigma}.
    \]
    This function is required to be smooth, and for every \(u\in\PreSectionSelectorBody{\sigma}\), its restriction to \(\KernelFiber{\sigma}{u}\) is required to be strictly convex and to attain a unique minimizer depending smoothly on \(u\); write
    \[
    \InducedKernelMinLift{\sigma}:\PreSectionSelectorBody{\sigma}\to\KernelBody{\sigma}.
    \]
    \item For each \(u\in\PreSectionSelectorBody{\sigma}\), define the selector-scale composed pre-pre-proto-kernel fiber
    \[
    \PrePreProtoKernelSelectorFiber{\sigma}{u}
    \equiv
    \left\{
    s\in\PrePreProtoKernelBody{\sigma}:
    \KernelProj{\sigma}\bigl(\PreKernelProj{\sigma}(\ProtoKernelProj{\sigma}(\PreProtoKernelProj{\sigma}(\PrePreProtoKernelProj{\sigma}(s)))))\bigr)=u
    \right\}.
    \]
    The restriction of \(\PrePreProtoKernelFunctional{\sigma}\) to that fiber is required to attain a unique minimizer depending smoothly on \(u\); write
    \[
    \PrePreProtoKernelSelectorMin{\sigma}:\PreSectionSelectorBody{\sigma}\to\PrePreProtoKernelBody{\sigma}.
    \]
    The restrictions of \(\InducedPreProtoKernelFunctional{\sigma}\) and \(\InducedProtoKernelFunctional{\sigma}\) to \(\PreProtoKernelSelectorFiber{\sigma}{u}\) and \(\ProtoKernelSelectorFiber{\sigma}{u}\), respectively, are also required to attain unique minimizers depending smoothly on \(u\); write
    \[
    \InducedPreProtoKernelSelectorMin{\sigma}:\PreSectionSelectorBody{\sigma}\to\PreProtoKernelBody{\sigma},
    \qquad
    \InducedProtoKernelSelectorMin{\sigma}:\PreSectionSelectorBody{\sigma}\to\ProtoKernelBody{\sigma}.
    \]
    Define the induced outer selector functional
    \[
    \InducedSelectorFunctional{\sigma}:\PreSectionSelectorSpace{\sigma}\to[0,\infty)
    \]
    by
    \[
    \InducedSelectorFunctional{\sigma}(u)
    \equiv
    \InducedKernelFunctional{\sigma}\bigl(\InducedKernelMinLift{\sigma}(u)\bigr)
    =
    \min_{w\in\KernelFiber{\sigma}{u}}
    \InducedKernelFunctional{\sigma}(w)
    \qquad
    \text{for every }
    u\in\PreSectionSelectorBody{\sigma}.
    \]
    This function is required to be smooth, and for every \(y\in\YoctoSelectorBody{\sigma}\), its restriction to \(\PreSectionSelectorFiber{\sigma}{y}\) is required to be strictly convex and to attain a unique minimizer depending smoothly on \(y\).
    \item For each \(y\in\YoctoSelectorBody{\sigma}\), the restriction of \(\InducedKernelFunctional{\sigma}\) to \(\KernelComposedFiber{\sigma}{y}\) is required to attain a unique minimizer depending smoothly on \(y\); write
    \[
    \InducedKernelComposedMin{\sigma}:\YoctoSelectorBody{\sigma}\to\KernelBody{\sigma}.
    \]
    Define the induced exact selector shell \(\InducedSelectorShell{\sigma}\), induced exact kernel shell \(\InducedKernelShell{\sigma}\), and induced exact pre-kernel shell \(\InducedPreKernelShell{\sigma}\) exactly as in Definition \ref{def:preprototarget}. They are required to satisfy the same smooth embeddedness and diffeomorphic projection conditions.
    \item For each \(y\in\YoctoSelectorBody{\sigma}\), the restrictions of \(\InducedProtoKernelFunctional{\sigma}\) and \(\InducedPreProtoKernelFunctional{\sigma}\) to \(\ProtoKernelFullFiber{\sigma}{y}\) and \(\PreProtoKernelFullFiber{\sigma}{y}\), respectively, are required to attain unique minimizers depending smoothly on \(y\); write
    \[
    \InducedProtoKernelFullMin{\sigma}:\YoctoSelectorBody{\sigma}\to\ProtoKernelBody{\sigma},
    \qquad
    \InducedPreProtoKernelFullMin{\sigma}:\YoctoSelectorBody{\sigma}\to\PreProtoKernelBody{\sigma}.
    \]
    Define the induced exact proto-kernel shell \(\InducedProtoKernelShell{\sigma}\) exactly as in Definition \ref{def:preprototarget}, and require
    \[
    \ProtoKernelProj{\sigma}\big|_{\InducedProtoKernelShell{\sigma}}
    :
    \InducedProtoKernelShell{\sigma}\xrightarrow{\cong}\InducedPreKernelShell{\sigma}
    \]
    to be a diffeomorphism.
    Define the induced exact pre-proto-kernel shell
    \[
    \InducedPreProtoKernelShell{\sigma}
    \equiv
    \left\{
    r\in\PreProtoKernelBody{\sigma}:
    \PreSectionSelectorProj{\sigma}\bigl(\KernelProj{\sigma}(\PreKernelProj{\sigma}(\ProtoKernelProj{\sigma}(\PreProtoKernelProj{\sigma}(r))))\bigr)\in\YoctoSelectorShell{\sigma},
    \;
    \InducedPreProtoKernelFunctional{\sigma}(r)
    =
    \min_{p\in\PreProtoKernelFullFiber{\sigma}{\PreSectionSelectorProj{\sigma}(\KernelProj{\sigma}(\PreKernelProj{\sigma}(\ProtoKernelProj{\sigma}(\PreProtoKernelProj{\sigma}(r)))))}}
    \InducedPreProtoKernelFunctional{\sigma}(p)
    \right\}.
    \]
    This shell is required to be a closed embedded submanifold, and
    \[
    \PreProtoKernelProj{\sigma}\big|_{\InducedPreProtoKernelShell{\sigma}}
    :
    \InducedPreProtoKernelShell{\sigma}\xrightarrow{\cong}\InducedProtoKernelShell{\sigma}
    \]
    is required to be a diffeomorphism.
    \item For each \(y\in\YoctoSelectorBody{\sigma}\), define the full pre-pre-proto-kernel fiber
    \[
    \PrePreProtoKernelFullFiber{\sigma}{y}
    \equiv
    \left\{
    s\in\PrePreProtoKernelBody{\sigma}:
    \PreSectionSelectorProj{\sigma}\bigl(\KernelProj{\sigma}(\PreKernelProj{\sigma}(\ProtoKernelProj{\sigma}(\PreProtoKernelProj{\sigma}(\PrePreProtoKernelProj{\sigma}(s))))))\bigr)=y
    \right\}.
    \]
    The restriction of \(\PrePreProtoKernelFunctional{\sigma}\) to that fiber is required to attain a unique minimizer depending smoothly on \(y\); write
    \[
    \PrePreProtoKernelFullMin{\sigma}:\YoctoSelectorBody{\sigma}\to\PrePreProtoKernelBody{\sigma}.
    \]
    The exact pre-pre-proto-kernel shell
    \[
    \PrePreProtoKernelShell{\sigma}
    \equiv
    \left\{
    s\in\PrePreProtoKernelBody{\sigma}:
    \PreSectionSelectorProj{\sigma}\bigl(\KernelProj{\sigma}(\PreKernelProj{\sigma}(\ProtoKernelProj{\sigma}(\PreProtoKernelProj{\sigma}(\PrePreProtoKernelProj{\sigma}(s))))))\bigr)\in\YoctoSelectorShell{\sigma},
    \;
    \PrePreProtoKernelFunctional{\sigma}(s)
    =
    \min_{t\in\PrePreProtoKernelFullFiber{\sigma}{\PreSectionSelectorProj{\sigma}(\KernelProj{\sigma}(\PreKernelProj{\sigma}(\ProtoKernelProj{\sigma}(\PreProtoKernelProj{\sigma}(\PrePreProtoKernelProj{\sigma}(s))))))}}
    \PrePreProtoKernelFunctional{\sigma}(t)
    \right\}
    \]
    is required to be a closed embedded submanifold, and
    \[
    \PrePreProtoKernelProj{\sigma}\big|_{\PrePreProtoKernelShell{\sigma}}
    :
    \PrePreProtoKernelShell{\sigma}\xrightarrow{\cong}\InducedPreProtoKernelShell{\sigma}
    \]
    is required to be a diffeomorphism.
    \item Benchmark faithfulness, meaning preservation of the benchmark outputs \eqref{eq:fixedoutputs}, no second operative metric, no enlarged low-energy matter law, and presentation as a deeper substrate-side source for the current pre-proto-kernel package rather than as a replacement benchmark theory.
\end{enumerate}
\end{definition}

\begin{remark}
Definition \ref{def:prepreproto} is intentionally written over the already fixed current pre-proto-kernel-body geometry. It does not reopen the benchmark selector base, the current selector-body geometry, the current kernel-body geometry, the current pre-kernel-body geometry, the current proto-kernel-body geometry, or the already-spent larger pre-selector, kernel, pre-kernel, and proto-kernel packages as the live burden. Its positive role is narrower: the current pre-proto-kernel functional, exact pre-proto-kernel shell, and current pre-proto-kernel minimizing data are generated as the effective multi-stage minimizing output of a deeper pre-pre-proto-kernel law.
\end{remark}

\begin{remark}
The label \emph{pre-pre-proto-kernel} is used here only as a local marker for a deeper layer below the current pre-proto-kernel frontier. It does not by itself assert a completed microscopic ontology or a final primitive law. Its role is only to name the next sharper transport package below the current pre-proto-kernel dynamics.
\end{remark}

\section{Recovery of the Current Pre-Selector Pre-Proto-Kernel Dynamics}

\begin{proposition}[The inner pre-pre-proto-kernel minimizer lift is well defined]
\label{prop:inner}
For every benchmark-faithful pre-pre-proto-kernel package, the minimizer lift \(\PrePreProtoKernelMinLift{\sigma}\) of Definition \ref{def:prepreproto} is a smooth section of \(\PrePreProtoKernelProj{\sigma}\), so that
\[
\PrePreProtoKernelProj{\sigma}\circ\PrePreProtoKernelMinLift{\sigma}
=
\mathrm{id}_{\PreProtoKernelBody{\sigma}}.
\]
In particular, the induced current pre-proto-kernel functional \(\InducedPreProtoKernelFunctional{\sigma}\) is a well-defined smooth function on \(\PreProtoKernelSpace{\sigma}\) whose values on \(\PreProtoKernelBody{\sigma}\) are the effective minimum values of the deeper pre-pre-proto-kernel functional on the inner pre-pre-proto-kernel fibers.
\end{proposition}

\begin{proof}
Definition \ref{def:prepreproto}(1) makes every inner pre-pre-proto-kernel fiber \(\PrePreProtoKernelFiber{\sigma}{r}\) nonempty, compact, and convex. Definition \ref{def:prepreproto}(2) states that \(\PrePreProtoKernelFunctional{\sigma}\) has a unique minimizer on each such fiber and that the minimizer depends smoothly on \(r\). Therefore \(\PrePreProtoKernelMinLift{\sigma}\) is a well-defined smooth map on \(\PreProtoKernelBody{\sigma}\). Each chosen minimizer lies in the fiber over the point at which it is selected, so
\[
\PrePreProtoKernelProj{\sigma}\bigl(\PrePreProtoKernelMinLift{\sigma}(r)\bigr)=r
\]
for every \(r\in\PreProtoKernelBody{\sigma}\). The remaining smoothness statement for \(\InducedPreProtoKernelFunctional{\sigma}\) is part of Definition \ref{def:prepreproto}(2), and the displayed formula identifies its values on \(\PreProtoKernelBody{\sigma}\) with the effective minimum values.
\end{proof}

\begin{proposition}[Current proto-kernel-scale pre-pre-proto-kernel minimization factors through the induced pre-proto-kernel minimization]
\label{prop:protofactor}
For every benchmark-faithful pre-pre-proto-kernel package and every \(z\in\ProtoKernelBody{\sigma}\),
\[
\PrePreProtoKernelProtoMin{\sigma}(z)
=
\PrePreProtoKernelMinLift{\sigma}\bigl(\InducedPreProtoKernelMinLift{\sigma}(z)\bigr),
\]
and hence
\[
\PrePreProtoKernelProj{\sigma}\bigl(\PrePreProtoKernelProtoMin{\sigma}(z)\bigr)
=
\InducedPreProtoKernelMinLift{\sigma}(z).
\]
\end{proposition}

\begin{proof}
Fix \(z\in\ProtoKernelBody{\sigma}\). Every \(s\in\PrePreProtoKernelProtoFiber{\sigma}{z}\) satisfies
\[
r\equiv\PrePreProtoKernelProj{\sigma}(s)\in\PreProtoKernelFiber{\sigma}{z}.
\]
By Proposition \ref{prop:inner},
\[
\PrePreProtoKernelFunctional{\sigma}(s)
\geq
\PrePreProtoKernelFunctional{\sigma}\bigl(\PrePreProtoKernelMinLift{\sigma}(r)\bigr)
=
\InducedPreProtoKernelFunctional{\sigma}(r).
\]
Definition \ref{def:prepreproto}(2) records that \(\InducedPreProtoKernelFunctional{\sigma}\) attains a unique minimum on \(\PreProtoKernelFiber{\sigma}{z}\), namely \(\InducedPreProtoKernelMinLift{\sigma}(z)\). Consequently, for every \(s\in\PrePreProtoKernelProtoFiber{\sigma}{z}\),
\[
\PrePreProtoKernelFunctional{\sigma}(s)
\geq
\InducedPreProtoKernelFunctional{\sigma}\bigl(\InducedPreProtoKernelMinLift{\sigma}(z)\bigr)
=
\PrePreProtoKernelFunctional{\sigma}\bigl(\PrePreProtoKernelMinLift{\sigma}(\InducedPreProtoKernelMinLift{\sigma}(z))\bigr).
\]
So \(\PrePreProtoKernelMinLift{\sigma}(\InducedPreProtoKernelMinLift{\sigma}(z))\) is a minimizer of \(\PrePreProtoKernelFunctional{\sigma}\) on \(\PrePreProtoKernelProtoFiber{\sigma}{z}\). Definition \ref{def:prepreproto}(3) says that minimizer is unique, hence the first displayed identity follows. The second identity is then Proposition \ref{prop:inner}.
\end{proof}

\begin{proposition}[The induced outer pre-proto-kernel dynamics form a benchmark-faithful package]
\label{prop:preproto}
For every benchmark-faithful pre-pre-proto-kernel package, the induced outer data
\[
\Bigl(
\PreProtoKernelSpace{\sigma},
\PreProtoKernelBody{\sigma},
\PreProtoKernelProj{\sigma},
\InducedPreProtoKernelFunctional{\sigma},
\InducedPreProtoKernelMinLift{\sigma},
\InducedPreProtoKernelPreKernelMin{\sigma},
\InducedPreProtoKernelKernelMin{\sigma},
\InducedPreProtoKernelSelectorMin{\sigma},
\InducedPreProtoKernelFullMin{\sigma},
\InducedPreProtoKernelShell{\sigma}
\Bigr)
\]
together with the induced current proto-kernel, pre-kernel, kernel, and selector functionals and shells of Definition \ref{def:prepreproto} form a benchmark-faithful proto-kernel-dynamics-generating substrate pre-selector pre-proto-kernel dynamics package on the recorded geometry in the sense of Definition \ref{def:preprototarget}.
\end{proposition}

\begin{proof}
Definition \ref{def:prepreproto}(2) states that \(\InducedPreProtoKernelFunctional{\sigma}\) is nonnegative and smooth, and that its restriction to each current pre-proto-kernel fiber \(\PreProtoKernelFiber{\sigma}{z}\) is strictly convex with a unique minimizer depending smoothly on \(z\). Definition \ref{def:prepreproto}(4) records the unique current pre-kernel-scale minimizer \(\InducedPreProtoKernelPreKernelMin{\sigma}\) on each current composed pre-proto-kernel fiber \(\PreProtoKernelPreKernelFiber{\sigma}{q}\) and defines the induced current pre-kernel functional with its unique minimizers on the current pre-kernel fibers. Definition \ref{def:prepreproto}(5) records the unique current-kernel-scale minimizer \(\InducedPreProtoKernelKernelMin{\sigma}\) on each current composed pre-proto-kernel fiber \(\PreProtoKernelKernelFiber{\sigma}{w}\), defines the induced current kernel functional, and supplies its unique minimizers on the current kernel fibers. Definition \ref{def:prepreproto}(6) records the unique selector-scale minimizer \(\InducedPreProtoKernelSelectorMin{\sigma}\) on each current composed pre-proto-kernel fiber \(\PreProtoKernelSelectorFiber{\sigma}{u}\), defines the induced selector functional, and keeps the selector, kernel, and pre-kernel shell structure fixed. Definition \ref{def:prepreproto}(7) records the unique full pre-proto-kernel minimizer \(\InducedPreProtoKernelFullMin{\sigma}\), the induced exact proto-kernel shell, the induced exact pre-proto-kernel shell, and the required diffeomorphic projections. Definition \ref{def:prepreproto}(8) gives benchmark faithfulness. These are precisely the requirements of Definition \ref{def:preprototarget}.
\end{proof}

\begin{proposition}[Current pre-kernel-scale pre-pre-proto-kernel minimization factors through the induced pre-proto-kernel pre-kernel minimization]
\label{prop:prekerfactor}
For every benchmark-faithful pre-pre-proto-kernel package and every \(q\in\PreKernelBody{\sigma}\),
\[
\PrePreProtoKernelPreKernelMin{\sigma}(q)
=
\PrePreProtoKernelMinLift{\sigma}\bigl(\InducedPreProtoKernelPreKernelMin{\sigma}(q)\bigr),
\]
and hence
\[
\PrePreProtoKernelProj{\sigma}\bigl(\PrePreProtoKernelPreKernelMin{\sigma}(q)\bigr)
=
\InducedPreProtoKernelPreKernelMin{\sigma}(q).
\]
\end{proposition}

\begin{proof}
Fix \(q\in\PreKernelBody{\sigma}\). Every \(s\in\PrePreProtoKernelPreKernelFiber{\sigma}{q}\) satisfies
\[
r\equiv\PrePreProtoKernelProj{\sigma}(s)\in\PreProtoKernelPreKernelFiber{\sigma}{q}.
\]
By Proposition \ref{prop:inner},
\[
\PrePreProtoKernelFunctional{\sigma}(s)
\geq
\PrePreProtoKernelFunctional{\sigma}\bigl(\PrePreProtoKernelMinLift{\sigma}(r)\bigr)
=
\InducedPreProtoKernelFunctional{\sigma}(r).
\]
Definition \ref{def:prepreproto}(4) records that \(\InducedPreProtoKernelFunctional{\sigma}\) attains a unique minimum on \(\PreProtoKernelPreKernelFiber{\sigma}{q}\), namely \(\InducedPreProtoKernelPreKernelMin{\sigma}(q)\). Consequently, for every \(s\in\PrePreProtoKernelPreKernelFiber{\sigma}{q}\),
\[
\PrePreProtoKernelFunctional{\sigma}(s)
\geq
\InducedPreProtoKernelFunctional{\sigma}\bigl(\InducedPreProtoKernelPreKernelMin{\sigma}(q)\bigr)
=
\PrePreProtoKernelFunctional{\sigma}\bigl(\PrePreProtoKernelMinLift{\sigma}(\InducedPreProtoKernelPreKernelMin{\sigma}(q))\bigr).
\]
So \(\PrePreProtoKernelMinLift{\sigma}(\InducedPreProtoKernelPreKernelMin{\sigma}(q))\) is a minimizer of \(\PrePreProtoKernelFunctional{\sigma}\) on \(\PrePreProtoKernelPreKernelFiber{\sigma}{q}\). Definition \ref{def:prepreproto}(4) says that minimizer is unique, hence the first displayed identity follows. The second identity is then Proposition \ref{prop:inner}.
\end{proof}

\begin{proposition}[Current-kernel-scale pre-pre-proto-kernel minimization factors through the induced pre-proto-kernel current-kernel minimization]
\label{prop:kernelfactor}
For every benchmark-faithful pre-pre-proto-kernel package and every \(w\in\KernelBody{\sigma}\),
\[
\PrePreProtoKernelKernelMin{\sigma}(w)
=
\PrePreProtoKernelMinLift{\sigma}\bigl(\InducedPreProtoKernelKernelMin{\sigma}(w)\bigr),
\]
and hence
\[
\PrePreProtoKernelProj{\sigma}\bigl(\PrePreProtoKernelKernelMin{\sigma}(w)\bigr)
=
\InducedPreProtoKernelKernelMin{\sigma}(w).
\]
\end{proposition}

\begin{proof}
Fix \(w\in\KernelBody{\sigma}\). Every \(s\in\PrePreProtoKernelKernelFiber{\sigma}{w}\) satisfies
\[
r\equiv\PrePreProtoKernelProj{\sigma}(s)\in\PreProtoKernelKernelFiber{\sigma}{w}.
\]
By Proposition \ref{prop:inner},
\[
\PrePreProtoKernelFunctional{\sigma}(s)
\geq
\PrePreProtoKernelFunctional{\sigma}\bigl(\PrePreProtoKernelMinLift{\sigma}(r)\bigr)
=
\InducedPreProtoKernelFunctional{\sigma}(r).
\]
Definition \ref{def:prepreproto}(5) records that \(\InducedPreProtoKernelFunctional{\sigma}\) attains a unique minimum on \(\PreProtoKernelKernelFiber{\sigma}{w}\), namely \(\InducedPreProtoKernelKernelMin{\sigma}(w)\). Consequently, for every \(s\in\PrePreProtoKernelKernelFiber{\sigma}{w}\),
\[
\PrePreProtoKernelFunctional{\sigma}(s)
\geq
\InducedPreProtoKernelFunctional{\sigma}\bigl(\InducedPreProtoKernelKernelMin{\sigma}(w)\bigr)
=
\PrePreProtoKernelFunctional{\sigma}\bigl(\PrePreProtoKernelMinLift{\sigma}(\InducedPreProtoKernelKernelMin{\sigma}(w))\bigr).
\]
So \(\PrePreProtoKernelMinLift{\sigma}(\InducedPreProtoKernelKernelMin{\sigma}(w))\) is a minimizer of \(\PrePreProtoKernelFunctional{\sigma}\) on \(\PrePreProtoKernelKernelFiber{\sigma}{w}\). Definition \ref{def:prepreproto}(5) says that minimizer is unique, hence the first displayed identity follows. The second identity is then Proposition \ref{prop:inner}.
\end{proof}

\begin{proposition}[Selector-scale pre-pre-proto-kernel minimization factors through the induced pre-proto-kernel selector minimization]
\label{prop:selectorfactor}
For every benchmark-faithful pre-pre-proto-kernel package and every \(u\in\PreSectionSelectorBody{\sigma}\),
\[
\PrePreProtoKernelSelectorMin{\sigma}(u)
=
\PrePreProtoKernelMinLift{\sigma}\bigl(\InducedPreProtoKernelSelectorMin{\sigma}(u)\bigr),
\]
and hence
\[
\PrePreProtoKernelProj{\sigma}\bigl(\PrePreProtoKernelSelectorMin{\sigma}(u)\bigr)
=
\InducedPreProtoKernelSelectorMin{\sigma}(u).
\]
\end{proposition}

\begin{proof}
Fix \(u\in\PreSectionSelectorBody{\sigma}\). Every \(s\in\PrePreProtoKernelSelectorFiber{\sigma}{u}\) satisfies
\[
r\equiv\PrePreProtoKernelProj{\sigma}(s)\in\PreProtoKernelSelectorFiber{\sigma}{u}.
\]
By Proposition \ref{prop:inner},
\[
\PrePreProtoKernelFunctional{\sigma}(s)
\geq
\PrePreProtoKernelFunctional{\sigma}\bigl(\PrePreProtoKernelMinLift{\sigma}(r)\bigr)
=
\InducedPreProtoKernelFunctional{\sigma}(r).
\]
Definition \ref{def:prepreproto}(6) records that \(\InducedPreProtoKernelFunctional{\sigma}\) attains a unique minimum on \(\PreProtoKernelSelectorFiber{\sigma}{u}\), namely \(\InducedPreProtoKernelSelectorMin{\sigma}(u)\). Consequently, for every \(s\in\PrePreProtoKernelSelectorFiber{\sigma}{u}\),
\[
\PrePreProtoKernelFunctional{\sigma}(s)
\geq
\InducedPreProtoKernelFunctional{\sigma}\bigl(\InducedPreProtoKernelSelectorMin{\sigma}(u)\bigr)
=
\PrePreProtoKernelFunctional{\sigma}\bigl(\PrePreProtoKernelMinLift{\sigma}(\InducedPreProtoKernelSelectorMin{\sigma}(u))\bigr).
\]
So \(\PrePreProtoKernelMinLift{\sigma}(\InducedPreProtoKernelSelectorMin{\sigma}(u))\) is a minimizer of \(\PrePreProtoKernelFunctional{\sigma}\) on \(\PrePreProtoKernelSelectorFiber{\sigma}{u}\). Definition \ref{def:prepreproto}(6) says that minimizer is unique, hence the first displayed identity follows. The second identity is then Proposition \ref{prop:inner}.
\end{proof}

\begin{proposition}[Benchmark-shell pre-pre-proto-kernel minimization factors through the induced pre-proto-kernel full minimization]
\label{prop:fullfactor}
For every benchmark-faithful pre-pre-proto-kernel package and every \(y\in\YoctoSelectorBody{\sigma}\),
\[
\PrePreProtoKernelFullMin{\sigma}(y)
=
\PrePreProtoKernelMinLift{\sigma}\bigl(\InducedPreProtoKernelFullMin{\sigma}(y)\bigr),
\]
and hence
\[
\PrePreProtoKernelProj{\sigma}\bigl(\PrePreProtoKernelFullMin{\sigma}(y)\bigr)
=
\InducedPreProtoKernelFullMin{\sigma}(y).
\]
\end{proposition}

\begin{proof}
Fix \(y\in\YoctoSelectorBody{\sigma}\). Every \(s\in\PrePreProtoKernelFullFiber{\sigma}{y}\) satisfies
\[
r\equiv\PrePreProtoKernelProj{\sigma}(s)\in\PreProtoKernelFullFiber{\sigma}{y}.
\]
By Proposition \ref{prop:inner},
\[
\PrePreProtoKernelFunctional{\sigma}(s)
\geq
\PrePreProtoKernelFunctional{\sigma}\bigl(\PrePreProtoKernelMinLift{\sigma}(r)\bigr)
=
\InducedPreProtoKernelFunctional{\sigma}(r).
\]
Definition \ref{def:prepreproto}(8) records that \(\InducedPreProtoKernelFunctional{\sigma}\) attains a unique minimum on \(\PreProtoKernelFullFiber{\sigma}{y}\), namely \(\InducedPreProtoKernelFullMin{\sigma}(y)\). Therefore, for every \(s\in\PrePreProtoKernelFullFiber{\sigma}{y}\),
\[
\PrePreProtoKernelFunctional{\sigma}(s)
\geq
\InducedPreProtoKernelFunctional{\sigma}\bigl(\InducedPreProtoKernelFullMin{\sigma}(y)\bigr)
=
\PrePreProtoKernelFunctional{\sigma}\bigl(\PrePreProtoKernelMinLift{\sigma}(\InducedPreProtoKernelFullMin{\sigma}(y))\bigr).
\]
So \(\PrePreProtoKernelMinLift{\sigma}(\InducedPreProtoKernelFullMin{\sigma}(y))\) is a minimizer of \(\PrePreProtoKernelFunctional{\sigma}\) on the full pre-pre-proto-kernel fiber \(\PrePreProtoKernelFullFiber{\sigma}{y}\). Definition \ref{def:prepreproto}(9) says that minimizer is unique, hence the first displayed identity follows. The second identity is then Proposition \ref{prop:inner}.
\end{proof}

\begin{proposition}[The exact pre-proto-kernel shell is the projected exact pre-pre-proto-kernel shell]
\label{prop:shell}
For every benchmark-faithful pre-pre-proto-kernel package,
\[
\InducedPreProtoKernelShell{\sigma}
=
\PrePreProtoKernelProj{\sigma}\bigl(\PrePreProtoKernelShell{\sigma}\bigr),
\]
the restriction
\[
\PrePreProtoKernelProj{\sigma}\big|_{\PrePreProtoKernelShell{\sigma}}
:
\PrePreProtoKernelShell{\sigma}\xrightarrow{\cong}\InducedPreProtoKernelShell{\sigma}
\]
is a diffeomorphism, and
\[
\PreProtoKernelProj{\sigma}\big|_{\InducedPreProtoKernelShell{\sigma}}
:
\InducedPreProtoKernelShell{\sigma}\xrightarrow{\cong}\InducedProtoKernelShell{\sigma}
\]
is a diffeomorphism.
\end{proposition}

\begin{proof}
Both diffeomorphisms are part of Definition \ref{def:prepreproto}(8)--(9). In particular, the first immediately implies
\[
\InducedPreProtoKernelShell{\sigma}
=
\PrePreProtoKernelProj{\sigma}\bigl(\PrePreProtoKernelShell{\sigma}\bigr).
\]
Nothing further is required.
\end{proof}

\begin{proposition}[No residual pre-proto-kernel-level freedom remains once the pre-pre-proto-kernel package is fixed]
\label{prop:nofreedom}
For every benchmark-faithful pre-pre-proto-kernel package, the current pre-proto-kernel functional \(\InducedPreProtoKernelFunctional{\sigma}\), the current pre-proto-kernel minimizer lift \(\InducedPreProtoKernelMinLift{\sigma}\), the current pre-proto-kernel pre-kernel minimizer \(\InducedPreProtoKernelPreKernelMin{\sigma}\), the current pre-proto-kernel current-kernel minimizer \(\InducedPreProtoKernelKernelMin{\sigma}\), the current pre-proto-kernel selector minimizer \(\InducedPreProtoKernelSelectorMin{\sigma}\), the current benchmark-shell-scale full pre-proto-kernel minimizer \(\InducedPreProtoKernelFullMin{\sigma}\), and the current exact pre-proto-kernel shell \(\InducedPreProtoKernelShell{\sigma}\) are uniquely determined by the deeper pre-pre-proto-kernel data. In particular, the current pre-proto-kernel package carries no residual benchmark-facing functional, shell, proto-kernel-scale minimizer-lift, pre-kernel-scale minimizer, current-kernel-scale minimizer, selector-scale minimizer, or benchmark-shell-scale minimizer freedom once the deeper pre-pre-proto-kernel package is fixed.
\end{proposition}

\begin{proof}
Proposition \ref{prop:inner} determines \(\InducedPreProtoKernelFunctional{\sigma}\) uniquely as the effective minimum value of \(\PrePreProtoKernelFunctional{\sigma}\) on the inner pre-pre-proto-kernel fibers. Definition \ref{def:prepreproto}(2) then uniquely determines \(\InducedPreProtoKernelMinLift{\sigma}\) as the unique minimizer of \(\InducedPreProtoKernelFunctional{\sigma}\) on each current pre-proto-kernel fiber. Definition \ref{def:prepreproto}(4) uniquely determines \(\InducedPreProtoKernelPreKernelMin{\sigma}\) on the current pre-kernel-scale composed pre-proto-kernel fibers, Definition \ref{def:prepreproto}(5) uniquely determines \(\InducedPreProtoKernelKernelMin{\sigma}\) on the current-kernel-scale composed pre-proto-kernel fibers, and Definition \ref{def:prepreproto}(6) uniquely determines \(\InducedPreProtoKernelSelectorMin{\sigma}\) on the selector-scale composed pre-proto-kernel fibers. Definition \ref{def:prepreproto}(8) uniquely determines \(\InducedPreProtoKernelFullMin{\sigma}\) and the exact pre-proto-kernel shell as the projected exact pre-pre-proto-kernel shell over \(\YoctoSelectorShell{\sigma}\). Therefore no benchmark-facing constituent at current pre-proto-kernel scope remains freely choosable once the deeper pre-pre-proto-kernel package is fixed.
\end{proof}

\begin{proposition}[All current downstream uses factor through the induced pre-proto-kernel dynamics]
\label{prop:downstream}
For every benchmark-faithful pre-pre-proto-kernel package, all current benchmark-facing downstream uses of the current pre-proto-kernel package on the active line factor through the induced pre-proto-kernel package of Proposition \ref{prop:preproto}. In particular, the current proto-kernel package \(\ProtoKernelPackage\), the current pre-kernel package \(\PreKernelPackage\), the current kernel package \(\KernelPackage\), the current pre-selector-minimizer-section-generating substrate selector dynamics package \(\PreSectionSelectorPackage\), the current observer-relevant pre-selector minimizer-section core \(\PreSelectorMinSectionCore\), the current selector package \(\YoctoSelectorPackage\), the current observer-relevant yocto-section minimizer-section core \(\YoctoMinSectionCore\), the current observer-relevant yocto-section minimizer-image core \(\YoctoImageCore\), the current observer-relevant zepto-section minimizer-section core \(\ZeptoMinSectionCore\), the current observer-relevant zepto-section minimizer-image core \(\ZeptoImageCore\), the current atto-section package \(\AttoSectionPackage\), the current femto-section package \(\FemtoSectionPackage\), the current pico-section package \(\PicoSectionPackage\), the current observer-relevant pico-section minimizer-section core \(\PicoMinSectionCore\), the current nano-section package \(\NanoSectionPackage\), the current micro-section package \(\MicroSectionPackage\), the current sub-section package \(\SubSectionPackage\), and the previously recorded shell-transport consequences follow unchanged.
\end{proposition}

\begin{proof}
By Proposition \ref{prop:preproto}, the induced outer data form a benchmark-faithful current pre-proto-kernel package of the exact type required by Definition \ref{def:preprototarget}. Definition \ref{def:downstream}(2) then records that, once such a package is available on the recorded geometry, the current proto-kernel package and all downstream benchmark-facing data and consequences follow unchanged. Therefore all current downstream uses already factor through the induced pre-proto-kernel package provided by the deeper pre-pre-proto-kernel dynamics.
\end{proof}

\begin{proposition}[The benchmark package remains fixed]
\label{prop:benchmark}
Every benchmark-faithful pre-proto-kernel-dynamics-generating substrate pre-selector pre-pre-proto-kernel dynamics package preserves the same benchmark one-metric package \eqref{eq:fixedoutputs}. In particular, the operative metric remains exactly \(\CompletedMetric\), the ordinary matter route is unchanged, there is no second operative metric, and the low-energy matter law is not enlarged.
\end{proposition}

\begin{proof}
This is part of benchmark faithfulness in Definition \ref{def:prepreproto}(10). The present note only derives the current pre-proto-kernel package from deeper substrate pre-pre-proto-kernel data. It does not alter the benchmark exact-closure package.
\end{proof}

\section{Main Theorem}

\begin{theorem}[Proto-kernel-dynamics-generating substrate pre-selector pre-proto-kernel dynamics from pre-proto-kernel-dynamics-generating substrate pre-selector pre-pre-proto-kernel dynamics]
\label{thm:main}
Fix the primitive continuation data of Definition \ref{def:fixed}, the recorded selector benchmark base of Definition \ref{def:selectorbase}, the recorded current selector-body geometry of Definition \ref{def:selectorgeometry}, the recorded current kernel-body geometry of Definition \ref{def:kernelgeometry}, the recorded current pre-kernel-body geometry of Definition \ref{def:prekernelgeometry}, the recorded current proto-kernel-body geometry of Definition \ref{def:protokernelgeometry}, the recorded current pre-proto-kernel-body geometry of Definition \ref{def:preprotogeometry}, the recorded downstream chain of Definition \ref{def:downstream}, and a benchmark-faithful pre-proto-kernel-dynamics-generating substrate pre-selector pre-pre-proto-kernel dynamics package in the sense of Definition \ref{def:prepreproto}. Then, for each sector \(\sigma\in\{\Car,\Cur,\Hom\}\), the current proto-kernel-dynamics-generating substrate pre-selector pre-proto-kernel dynamics are no longer an unexplained primitive stopping point on the active line: they are recovered canonically from the deeper pre-pre-proto-kernel dynamics. More precisely:
\begin{enumerate}
    \item the current pre-proto-kernel functional is induced as the effective minimum value of the deeper pre-pre-proto-kernel functional on the inner pre-pre-proto-kernel fibers by Proposition \ref{prop:inner};
    \item the current pre-proto-kernel minimizer lift on the current pre-proto-kernel fibers over \(\ProtoKernelBody{\sigma}\) is the outer projection of the unique current proto-kernel-scale pre-pre-proto-kernel minimizer by Proposition \ref{prop:protofactor};
    \item the induced outer data form a benchmark-faithful current pre-proto-kernel package on the recorded geometry by Proposition \ref{prop:preproto};
    \item the current pre-proto-kernel pre-kernel minimizer is the outer projection of the unique current pre-kernel-scale pre-pre-proto-kernel minimizer by Proposition \ref{prop:prekerfactor};
    \item the current pre-proto-kernel current-kernel minimizer is the outer projection of the unique current-kernel-scale pre-pre-proto-kernel minimizer by Proposition \ref{prop:kernelfactor};
    \item the current pre-proto-kernel selector minimizer is the outer projection of the unique selector-scale pre-pre-proto-kernel minimizer by Proposition \ref{prop:selectorfactor};
    \item the current benchmark-shell-scale full pre-proto-kernel minimizer is the outer projection of the unique full pre-pre-proto-kernel minimizer by Proposition \ref{prop:fullfactor};
    \item the current exact pre-proto-kernel shell is exactly the projected exact pre-pre-proto-kernel shell, and it still projects diffeomorphically to the induced exact proto-kernel shell, by Proposition \ref{prop:shell};
    \item the current pre-proto-kernel package carries no residual benchmark-facing functional, shell, proto-kernel-scale minimizer-lift, pre-kernel-scale minimizer, current-kernel-scale minimizer, selector-scale minimizer, or benchmark-shell-scale minimizer freedom once the deeper pre-pre-proto-kernel package is fixed by Proposition \ref{prop:nofreedom};
    \item all current downstream benchmark-facing uses of the pre-proto-kernel package already factor through the induced outer package by Proposition \ref{prop:downstream};
    \item the benchmark boundary remains fixed by Proposition \ref{prop:benchmark};
    \item consequently the remaining honest burden below the current pre-selector pre-proto-kernel frontier is no longer the proto-kernel-dynamics-generating substrate pre-selector pre-proto-kernel dynamics \(\PreProtoKernelPackage\) as such. What remains is to derive or justify the deeper pre-proto-kernel-dynamics-generating substrate pre-selector pre-pre-proto-kernel dynamics \(\PrePreProtoKernelPackage\) from still more primitive substrate structure, or else prove a still sharper theorem-level collapse strictly inside that deeper package.
\end{enumerate}
\end{theorem}

\begin{proof}
Item (1) is Proposition \ref{prop:inner}. Item (2) is Proposition \ref{prop:protofactor}. Item (3) is Proposition \ref{prop:preproto}. Item (4) is Proposition \ref{prop:prekerfactor}. Item (5) is Proposition \ref{prop:kernelfactor}. Item (6) is Proposition \ref{prop:selectorfactor}. Item (7) is Proposition \ref{prop:fullfactor}. Item (8) is Proposition \ref{prop:shell}. Item (9) is Proposition \ref{prop:nofreedom}. Item (10) is Proposition \ref{prop:downstream}. Item (11) is Proposition \ref{prop:benchmark}. Item (12) is the routing consequence of the previous items together with the active safeguard against counting another same-output deeper-origin relay as new front-line progress when the live benchmark-relevant pre-proto-kernel package is unchanged.
\end{proof}

\begin{corollary}[What must be done next]
\label{cor:next}
After Theorem \ref{thm:main}, the next honest benchmark-facing manuscript must either:
\begin{enumerate}
    \item derive or justify the pre-proto-kernel-dynamics-generating substrate pre-selector pre-pre-proto-kernel dynamics \(\PrePreProtoKernelPackage\) from still more primitive substrate structure in a way that changes benchmark-facing status;
    \item identify a genuinely smaller theorem-level observer-relevant datum strictly inside \(\PrePreProtoKernelPackage\) and prove a sharper collapse to it;
    \item do either of the previous two while preserving the same benchmark one-metric package and the same claim boundary.
\end{enumerate}
Another same-output deeper-origin relay that merely preserves the current pre-pre-proto-kernel dynamics \(\PrePreProtoKernelPackage\) and therefore merely reproduces the same current pre-proto-kernel package \(\PreProtoKernelPackage\), the same current proto-kernel package \(\ProtoKernelPackage\), the same current pre-kernel package \(\PreKernelPackage\), the same current kernel package \(\KernelPackage\), the same current selector-dynamics package \(\PreSectionSelectorPackage\), the same current observer-relevant pre-selector minimizer-section core \(\PreSelectorMinSectionCore\), the same current selector package \(\YoctoSelectorPackage\), the same current observer-relevant yocto-section minimizer-section core \(\YoctoMinSectionCore\), the same current observer-relevant yocto-section minimizer-image core \(\YoctoImageCore\), and the same downstream observer-relevant zepto-section, atto-section, femto-section, pico-section, nano-section, micro-section, and sub-section consequences, another re-expansion back to the current pre-proto-kernel package, the current proto-kernel package, the current pre-kernel package, the current kernel package, the current selector-dynamics package, the larger pre-selector package, or older yocto-section presentations as if they were still independently live, or any move that introduces a second operative metric, enlarges the ordinary matter law, or uses stronger promotion language does not satisfy the present burden.
\end{corollary}

\begin{proof}
Theorem \ref{thm:main} shows that the current pre-proto-kernel package is no longer the deepest unexplained datum on the active line. Therefore a subsequent manuscript must improve on the deeper pre-pre-proto-kernel dynamics rather than merely restate the already induced pre-proto-kernel package.
\end{proof}

\section{Conclusion}

The positive result of this note is that the current proto-kernel-dynamics-generating substrate pre-selector pre-proto-kernel dynamics are no longer primitive on the active derivational line. Once one works with pre-proto-kernel-dynamics-generating substrate pre-selector pre-pre-proto-kernel dynamics over the already fixed current pre-proto-kernel-body geometry, the current pre-proto-kernel functional is forced as the effective minimum value of the deeper pre-pre-proto-kernel law, the current pre-proto-kernel minimizer lift is forced as the outer projection of the unique current proto-kernel-scale pre-pre-proto-kernel minimizer, the current pre-proto-kernel pre-kernel minimizer is forced as the outer projection of the unique current pre-kernel-scale pre-pre-proto-kernel minimizer, the current pre-proto-kernel current-kernel minimizer is forced as the outer projection of the unique current-kernel-scale pre-pre-proto-kernel minimizer, the current pre-proto-kernel selector minimizer is forced as the outer projection of the unique selector-scale pre-pre-proto-kernel minimizer, the current benchmark-shell-scale full pre-proto-kernel minimizer is forced as the outer projection of the unique full pre-pre-proto-kernel minimizer, and the current exact pre-proto-kernel shell is forced as the projected exact pre-pre-proto-kernel shell over the recorded selector benchmark shell. All current downstream benchmark-facing uses then factor through that induced pre-proto-kernel package together with the already fixed selector benchmark base, selector-body geometry, kernel-body geometry, pre-kernel-body geometry, and proto-kernel-body geometry.

This preserves the same benchmark one-metric package and the same claim boundary. The result does not provide a completed first-principles derivation of GR from substrate microphysics, and it does not license stronger promotion language. Its narrower benchmark-facing gain is that the remaining honest burden now sits one level lower again: not on the current pre-selector pre-proto-kernel dynamics themselves, but on the deeper pre-selector pre-pre-proto-kernel dynamics that canonically generate them, or on a still smaller collapse strictly inside that deeper package.

\end{document}
